\section{Reproducibility statement}

% It is important that the work published in ICLR is reproducible. Authors are strongly encouraged to include a paragraph-long Reproducibility Statement at the end of the main text (before references) to discuss the efforts that have been made to ensure reproducibility. This paragraph should not itself describe details needed for reproducing the results, but rather reference the parts of the main paper, appendix, and supplemental materials that will help with reproducibility. For example, for novel models or algorithms, a link to an anonymous downloadable source code can be submitted as supplementary materials; for theoretical results, clear explanations of any assumptions and a complete proof of the claims can be included in the appendix; for any datasets used in the experiments, a complete description of the data processing steps can be provided in the supplementary materials. Each of the above are examples of things that can be referenced in the reproducibility statement. This optional reproducibility statement is not part of the main text and therefore will not count toward the page limit. 

We provide the code for our data generation pipeline, along with detailed instructions for executing the pipeline end-to-end, as well as sample dataset files in the supplementary materials. The main paper and appendix further document key implementation details, including prompt templates, hyperparameter configurations used during fine-tuning, and extensions of our data analysis and fine-tuning experiments. After publication, we plan to release the full codebase in a public GitHub repository and make our datasets publicly available on the HuggingFace platform.